% Folder 144, batch 10 (archivists' pages 181-183).
% Pass: Opus 5.5 (claude-opus-5-5), 2026-10-03 — first pass, unchecked against the pages by a human.
%
% Page 183 of the batch is a typed examination paper (DEUG A, juin 1981) with
% nothing in his hand: skipped. The pencilled numbers on the leaves read
% 182, 183, 184; \page{} follows the batch numbering 181-183 given by the
% batch split (page 1 of the batch = archivists' page 181).
\documentclass[11pt,a4paper]{article}
\input{../preamble/grothendieck}

\folder{144}
\batch{10}
\pages{181}{183}
\foldertitle{[Suite autour de Teichmüller dont] Teichmülleries, 1981-1982 : notes manuscrites (1981-1983, s.d.), lettres (1981, s.d.).}
\dating{1981-1983}
\watermark{Édition de démonstration}

\begin{document}

\page{181}
\note{la page s'ouvre sur une liste numérotée de conditions sur une courbe $X$ ; ce qui précède la condition 1° n'est pas dans ce lot.}

\marginal{$g$, $\mu$}

\drawing{dans la marge de gauche, plusieurs croquis de courbes à points doubles ordinaires : des branches qui se croisent, une boucle se recoupant elle-même, et un cycle de six droites se coupant deux à deux en hexagone.}

1°) \add{$X$ propre,} connexe, lisse \add{\uncertain{sauf}} pts doubles ordinaires \add{(diviseur $D$)}, dim 1

2°) $\chi(X, \mathcal{O}_X) = 1-g$
\note{la formule est surchargée : le second argument de $\chi$ est récrit par-dessus un premier jet ; la lecture $\mathcal{O}_X$ est probable sans être sûre.}

i.e.\ si $\widetilde{X}$ la normalisée, $g_i$ les genres, $\mu$ le nb des pts doubles
\note{un petit signe en exposant après $\widetilde{X}$ reste \ill{}.}
\[
  g-1 = \sum (g_i - 1) + \mu
\]
\[
  g = \sum g_i + \mu + 1 - s
\]
$s$ nb des comp.\ irr., $\mu$ nb des pts doubles (NB on a toujours $\mu + 1 - s \geq 0$ à cause connexité)

3°) $S \subset X$ \struck{\uncertain{multisection} \ill{} \uncertain{diviseur}} \add{sous-schéma fini ét., de rang $\nu$} contenu dans l'\uncertain{ens} des pts de lissité de $X$

\marginal{« \uncertain{types} d'\ill{} \uncertain{inf.}\ de la \uncertain{rigidité} »}\note{cette note marginale est reliée par une accolade aux conditions 4° et 5°.}

4°) Si $X_i$ comp.\ irréd.\ \add{avec} \struck{de genre} $\widetilde{X}_i$ genre 1, alors $X_i$ contient un pt double \uncertain{ou} $X_i \cap S \neq \emptyset$, i.e.\ $X_i \cap (D \cup S) \neq \emptyset$ ($\Longleftrightarrow$ \struck{$X$} \uncertain{sauf si} $X$ lisse de genre 1, $S = \emptyset$)

5°) Si $X_i$ comp.\ irréd.\ avec $\widetilde{X}_i$ de genre 0, alors l'image inverse de $(D \cup S)$ dans $\widetilde{X}_i$ a au moins trois pts.

\begin{tikzcd}
\widetilde{D} \arrow[r, hook] \arrow[dr, "\text{surj}"'] & \widetilde{D} \amalg S \arrow[r, "\text{surj}"] & \pi_0(\widetilde{X}) \arrow[r, "g"] & \mathbb{N} \\
 & D & &
\end{tikzcd}

$\pi_0(\widetilde{X}) = \Delta$ ; $\widetilde{D} \to D$ fibres de card $\leq 2$
\note{sous $\widetilde{D} \amalg S$, une accolade porte un signe \ill{} suivi d'un \uncertain{$\widetilde{\Delta}$}.}

À droite, sous une accolade :
\[
  \operatorname{card} S = \nu, \qquad \sum_{\delta \in \Delta} (g_\delta - 1) + \underbrace{\mu}_{\operatorname{card} \widetilde{D}} = g - 1,
\]
\[
  \operatorname{card}\bigl(D \amalg_{\widetilde{D}} \Delta\bigr) = 1
\]
\struck{fibres de $\widetilde{D} \to \Delta$ de card $\leq 2$}
\note{$\nu$ et $g$ sont entourés. Au-dessus du premier $D$ de la dernière égalité, un « \uncertain{$2\mu$} » paraît biffé ; sous $\mu$ il écrit « card $\widetilde{D}$ », alors que $\operatorname{card}\widetilde{D} = 2\mu$ (cf.\ p. 182, où $\operatorname{card}\widetilde{D} = 2\mu$).}

\page{182}
\uncertain{Squelette} de $X$ (sous hyp.\ \uncertain{dessus})
\begin{align*}
  \widetilde{D} &= \text{image inv.\ de } D \text{ dans } \widetilde{X} \\
  \widetilde{S} &= \text{image inv.\ de } S \text{ dans } \widetilde{X} \qquad (\text{NB } \widetilde{S} \xrightarrow{\sim} S) \\
  \Delta &= \pi_0(\widetilde{X}) \simeq \text{ens des comp.\ irréd.\ de } X \\
  D &= (\text{pour mémoire}) \text{ ens des pts doubles}
\end{align*}
\note{à la deuxième ligne, « image inv. de » est rendu par un trait de répétition.}

$\Sigma$ :

\begin{tikzcd}
\widetilde{D} \arrow[r, hook, "\text{can}"] \arrow[dr, "q\ (\text{surj})"'] & \widetilde{D} \amalg \widetilde{S} \arrow[r, "p\ (\text{surj})"] & \Delta \arrow[r, "g"] & \mathbb{N} \\
 & D & &
\end{tikzcd}

a) $D \amalg_{\widetilde{D}} \Delta$ de card.\ 1

b) \struck{les fibres de $q : \widetilde{D} \to D$ de card $\leq 2$} \add{de degré 2}

c) $\displaystyle \sum_{\delta \in \Delta} \bigl(g(\delta) - 1\bigr) + \underbrace{\operatorname{card} D}_{\mu} = g - 1$
\note{$g$ est entouré. Au-dessus du $D$ de « card $D$ », une marque, peut-être un tilde biffé ; l'accolade donne $\mu$, qui est $\operatorname{card} D$.}

d) $\operatorname{card} \widetilde{S} = \nu$ \note{$\nu$ entouré.}

e) Si $\delta \in \Delta$ tel que $g(\delta) = 1$, alors $p^{-1}(\delta) \neq \emptyset$ ($\Longleftrightarrow$ on n'a pas $\widetilde{D} = \widetilde{S} = \emptyset$, $\Delta = \{\delta\}$, $g(\delta) = 1$)

f) Si $\delta \in \Delta$ tel que $g(\delta) = 0$, alors $\operatorname{card} p^{-1}(\delta) \geq 3$

\[
  \mu(\Sigma) = \sum_{\delta \in \Delta} \mu(\delta, \Sigma)
\]
où $\mu(\delta, \Sigma)$ ne dépend que de $g(\delta) \overset{\text{déf}}{=} g_\delta$, et de $\operatorname{card}\bigl(p^{-1}(\delta)\bigr) \overset{\text{déf}}{=} \nu_\delta$,
\[
  \mu(\delta, \Sigma) = \mu(g_\delta, \nu_\delta) \overset{\text{déf}}{=} 3(g_\delta - 1) + \nu_\delta
\]
\note{avant $\mu(g_\delta, \nu_\delta)$, quelques signes biffés \ill{}. La définition était d'abord une accolade à deux cas : à la suite de $3(g_\delta - 1) + \nu_\delta$ un « si $g_\delta \geq 1$ » est barré, et la seconde ligne de l'accolade (\ill{}, « si $g_\delta$ … 1 ») est entièrement biffée.}
\begin{align*}
  \mu(\Sigma) &= \sum_{\delta} \mu(\delta, \Sigma) = 3 \sum (g_\delta - 1) + \underbrace{\sum \nu_\delta}_{\operatorname{card}\widetilde{D} + \operatorname{card}\widetilde{S}} \\
  &= 3 \sum (g_\delta - 1) + 2\mu + \nu \\
  &= 3 \underbrace{\Bigl(\sum_{\delta} (g_\delta - 1) + \mu\Bigr)}_{g-1} + \nu - \mu \\
  &= \underbrace{\bigl(3(g-1) + \nu\bigr)}_{\mu(g, \nu)} - \mu = \mu(g, \nu) - \mu
\end{align*}
\note{sur la page, $\operatorname{card}\widetilde{D}$ et $\operatorname{card}\widetilde{S}$ portent eux-mêmes des accolades $2\mu$ et $\nu$. Le $\mu(\Sigma)$ de tête est récrit par-dessus un premier jet. Le calcul s'arrête en bas de la page ; la page suivante du lot ne le continue pas, et ce qui suit éventuellement n'est pas dans ce lot.}

\end{document}
